\section{Institutional setting and data}
\subsection{Index construction and event timing}

OMXSB is designed to represent a selection of large and actively traded securities listed in Stockholm. The current methodology combines turnover ranking with free-float capitalization within industry supersectors. High-turnover securities qualify for inclusion, while the least liquid portion of the universe is ineligible; the industry selection targets at least 85\% of each supersector's free-float capitalization, subject to the eligibility rules. The index is free-float capitalization weighted. This combination of rules makes membership related to firm size, past trading and industry composition \citep{nasdaq2026}.

The current review calendar uses reference dates at the end of April and October, announcements after the close on the sixth trading day before effectiveness, and implementation at the opening of the first trading day in June or December. Historical event dates are taken from the event-specific source list, not retroactively assigned from this current calendar. The sample includes earlier January and July implementations. The earlier rules document provides a historical institutional reference, but neither that document nor the current methodology substitutes for a notice-by-notice event calendar \citep{nasdaq2013,nasdaq2026}.

Selection rules also change within the sample. The May 2012 notice explicitly announces revised industry-diversification criteria based on the Industry Classification Benchmark, alongside 11 additions and 24 removals \citep{nasdaq2012review}. Historical period comparisons therefore span institutional changes as well as changes in investor demand and market conditions. The unusually large review should not be treated as evidence that the assignment process was constant over time. This documented change motivates sensitivity to review influence; it does not supply an exogenous break design.

We denote the recorded announcement trading date by AD and the effective trading date by ED. If implementation occurs at the opening of ED, the previous trading day is the final trading day under the old composition. This distinction is important for the volume analysis. A large volume share at ED-1 is consistent with portfolio adjustment before effectiveness, but daily volume does not identify auction volume. Announcement timestamps also matter. The symmetric three-day announcement window includes the next trading day when information is released after the close; daily data still cannot isolate a narrow announcement-time response.

Membership selection need not be independent of future returns. Past price increases can raise capitalization, past trading can affect liquidity eligibility, and corporate developments can influence both eligibility and expected cash flows. A predictable rule therefore does not by itself make membership changes as good as random. The economic interpretation below explicitly allows selection and information effects to accompany trading demand.

\subsection{Sample construction and missing observations}

The frozen event file contains 284 additions and 252 deletions. Each row includes the company, security identifiers, event type, review cycle, announcement and effective dates, and a source URL. The imported price-event panel contains 487 candidates. Its reduction from 536 reflects the available historical security mappings and return data, rather than an economic exclusion of all missing firms. Missing observations are not assumed to be random.

\begin{table}[!htbp]
\centering
\caption{Construction of the return sample}\label{tab:sample}
\small
\setlength{\tabcolsep}{5pt}
\begin{tabular}{lrr}
\toprule
Stage & Additions & Deletions \\
\midrule
Documented membership changes & 284 & 252 \\
Imported return candidates & 261 & 226 \\
After deletion exclusions & 261 & 216 \\
Complete main return window & 256 & 215 \\
\bottomrule
\end{tabular}
\begin{minipage}{\linewidth}\footnotesize\vspace{5pt}
The source list contains 536 membership changes in 39 review cycles from 2007 through the first review of 2026. The imported return panel contains 487 candidates. Ten deletions flagged for offers or delisting are excluded; six remaining candidates lack a complete main window. The final sample contains 471 event observations. Each event type spans 38 review cycles; their union spans 39. Source: frozen event files and independent reconstruction, audit/results/sample\_flow.csv.
\end{minipage}
\end{table}


The main deletion sample excludes ten candidates flagged for an offer, merger-related event, or delisting under the inherited screening rule. Other deletions flagged only for changes in share counts remain in the return analysis but may be excluded from holdings-based calculations. Requiring a complete main return window leaves 256 additions and 215 deletions. The total sample spans 39 review cycles, although each event type is observed in only 38 of them. The asymmetry is why type-specific cluster counts must be reported rather than inferred from the pooled sample.

We divide events by effective year into 2007--2011, 2012--2016, 2017--2021 and 2022--2026. These are calendar blocks of equal nominal length, not equally informative samples: the study begins in July 2007 and stops after the June 2026 review. Event counts, coverage and review sizes also differ. Pooling the first two and last two blocks produces the earlier and later decades. The period labels should not be read as evidence that the endpoints were independently selected before the research process. The repository records earlier period definitions and subsequent revisions.

Selection has two distinct consequences. It can change the mean return when omitted securities differ systematically from included securities. It can also change the precision of that mean through the size and balance of review clusters. Later complete return samples are larger, but their availability does not repair missing historical observations. We report sample counts at every material step and avoid inferring an economic time trend directly from the growth in usable observations.

\subsection{Returns, prices and trading data}

Daily security data are supplied through LSEG. The principal return file includes closing prices, daily total returns, volume, market capitalization, shares outstanding and closing bid and ask quotes. Separate index files supply OMXSBGI and other market series. Broader stock panels support the liquidity and comparison exercises. The data are licensed research inputs; a public replication package can identify them and document transformations without redistributing the source observations.

The inherited daily return construction uses the reported total return when available and falls back to the percentage change in the closing price otherwise. It requires every abnormal-return observation within an event window to be present. The independent audit reproduces the announcement, implementation and total-window CARs for all 487 imported candidates. Seven candidates use at least one price-return fallback within the main window, covering 104 event-day observations. Eliminating implicit forward filling in that fallback does not change a main-window CAR in this frozen sample. This check validates the implemented arithmetic; it does not establish that every vendor observation or historical identifier is correct.

OMXSBGI is a gross total-return benchmark, so the benchmark and reported stock total returns are aligned more closely than a price-only index would permit. Nevertheless, a market-adjusted return imposes a unit market exposure and zero intercept. Differences in risk exposure, sector performance, dividends in fallback observations and changes in benchmark membership can remain in the residual. We therefore call the outcome market-adjusted CAR and retain the benchmark as part of its definition.

\subsection{Fund holdings and the demand sample}

The holdings analysis uses LSEG/Lipper fund positions around the events and a dated classification assembled from fund mandates and source documents. Three nested groups are retained: B contains strict OMXSB followers; V adds benchmark variants; C adds relevant mixed-fund exposure. The broad C measure is the main descriptive demand proxy, while B and V reveal how the measurement changes with the investor definition. An index-trade diagnostic checks the frequency of purchases of additions and disposals of deletions. The production code reads a dated mandate list; it does not mechanically reclassify every fund-year from that diagnostic alone.

Eligible fund-event comparisons require a report within 120 calendar days before AD and a report from ED through 45 calendar days afterwards, together with a stable classification across the two reporting dates. The code avoids treating a stale pre-event holding returned as an ``after'' observation as a new position. It nevertheless infers some zero positions from the presence of a report elsewhere in the observed fund data. That is a coverage assumption: the absence of a security row need not always mean a true zero if an underlying report or vendor extract is incomplete.

The change in reported shares held is normalized by the event's shares-outstanding denominator and sign-adjusted for deletions. Events with flagged splits are excluded from the principal ratio estimates to reduce inconsistencies between adjusted holdings and outstanding shares. The measurement still requires historical identifiers and share-class denominators to be correct. The longer interval between holdings reports also contains trading unrelated to the index event. Accordingly, the holdings change is not an observed implementation-day order or an independently assigned demand shock.

The resulting sample is much smaller than the return sample. The first addition block has eight usable C observations across five reviews, while the final block has 78 across nine reviews. The later decade contains 136 usable additions and 108 deletions. Comparing the full-sample CAR with demand measured only for those subsets would confound differences in price response with differences in coverage. Every reported ratio instead uses the same events for its numerator and denominator.

\subsection{Other measurements and provenance}

The study also uses MSCI membership-change records, a targeted company-news screen, a set of size-screened comparison securities, and share-class pairs. The MSCI exercise asks whether removing overlapping index events changes the OMXSB return pattern. The news screen covers the 24 most extreme inherited returns. It is not a comprehensive news audit and cannot certify the remaining sample as free of firm information. The share-class data include supplemental price sources, which motivates a separate analysis excluding source changes.

All reported estimates refer to one frozen repository version. Forty-three input files pass their recorded checksum tests. The archived code is run separately from the frozen output, and the independent return reconstruction does not import its return or inference functions. Source URLs and retrieval outcomes are preserved in the audit package. Complete notice reconciliation and report-level holdings checks remain unresolved; the paper does not recast an inherited verification flag as a new independent source verification.
